\section{Introduction}
Since the emergence of the Diffusion Transformer-based (DiT) video diffusion models, the field of general video generation, including Text-to-Video and Image-to-Video~\cite{bar2024lumiere,svd,ayl,guo2023animatediff,zhou2022magicvideo,walt,wang2023modelscope,vdm,videogan,cvideogan,singer2022make,text2video,villegas2022phenaki,lin2025apt} has made significant progress in producing highly realistic video content. A key factor driving this advancement is the large-scale training data, typically formatted as video-text pairs. Expanding the training dataset enables DiT networks to learn motion priors for various objects and scenes, resulting in strong generalization capabilities during inference.

Building upon these pretrained video diffusion networks, end-to-end human animation models, especially for audio-driven talking human generation, have developed rapidly since last year~\cite{he2023gaia,tian2024emo,xu2024hallo,wang2024vexpress,chen2024echomimic,xu2024vasa,stypulkowski2024diffused,jiang2024loopy,lin2025cyberhost}. 
Despite achieving realistic results, these models are trained on highly filtered datasets to simplify the learning process, restricting their applicability to limited scenarios. 
For instance, most existing end-to-end audio-driven models are limited to facial or portrait images captured from a front-facing perspective with a static background. To date, no prior work has attempted to scale up training data for more generalizable human animation.

Scaling up human animation data may seem straightforward, but unfortunately it is not. Directly adding more data is not always beneficial for network training. 
Take audio-conditioned models as an example: audio is primarily associated with facial expressions and has little correlation with body poses, background motion, camera movement, or lighting changes. As a result, raw training data must be filtered and cropped to minimize the influence of these unrelated factors. Additionally, audio-conditioned models often undergo further data cleaning based on lip-sync accuracy, which is also important to stabilize training. 
% Similarly, pose-conditioned models require extensive filtering, cropping, and cleaning. 
Unfortunately, these processes discard a substantial amount of data, making dataset scaling a futile effort or difficult to achieve, despite much of the discarded data containing valuable motion patterns essential for training data expansion. For example, recent state-of-the-art methods mention that after rigorous cleaning, less than 10\% of the data is retained, rendering direct data scaling highly cost-ineffective.

In this paper, we address the challenges of scaling up audio-driven human animation data and models. Our key insight is that incorporating multiple conditioning signals beyond audios, such as text and pose, during training can significantly reduce data wastage. This approach offers two main advantages. On one hand, data that would otherwise be discarded for single-condition models (e.g., audio-driven) can be leveraged in tasks with weaker or more general conditions, such as text. Training on such data allows the model to learn more diverse motion patterns, mitigating the limitations imposed by data filtering. On the other hand, different conditioning signals can complement each other. For example, while audio alone cannot precisely control body poses, stronger conditions such as body poses can provide additional guidance. By integrating stronger conditioning signals alongside audio data, we aim to reduce overfitting and improve the generalization of generated results.


Building on the above considerations, we designed an omni-conditions training strategy, which follows two key training principles: (1) tasks with stronger conditioning can leverage those with weaker conditioning, along with their associated data, to scale up training, and (2) the stronger the conditioning signal, the lower its training ratio should be. To implement these strategies, we develop OmniHuman, a mixed conditioned human video generation model based on the advanced video generation architecture, DiT~\cite{dit,sd3}. Although OmniHuman mainly focuses on audio-driven human video generation, it is trained with three motion-related conditions, including text, audio, and pose, ranging from weak to strong. This approach effectively addresses the challenge of data scaling in end-to-end audio-driven frameworks, allowing the model to learn natural motion patterns from large-scale data and support various input forms.



In summary, we introduce OmniHuman, an audio-driven human video generation model that leverages our omni-conditions training strategy to integrate various motion-related conditions and their corresponding data. Unlike existing methods that reduce data due to strict filtering, our approach benefits from large-scale mixed-conditioned data, enabling OmniHuman to generate highly realistic and expressive human motion videos given a reference image and the driving audio. It adapts to various portrait types and aspect ratios, and significantly improves gesture generation, which is a longstanding challenge for previous methods. Additionally, it accommodates diverse image styles and background contents. To the best of our knowledge, OmniHuman is the first solution capable of audio-driven human video generation on input images with any body proportions and image styles, and also supports auxiliary pose driving.
